Se evaluaron modelos de primer y segundo orden con retardo para describir
el comportamiento térmico de un horno eléctrico. Los modelos generados en
Matlab fueron los que mejor se ajustaron a los datos experimentales.

Los métodos de primer orden, en particular el obtenido por promedio de
puntos, presentaron menor error que los métodos de segundo orden.

Finalmente, se encontró que el método de Stark presentó inconsistencias en
la estimación del retardo, lo que afectó su precisión como predictor.
